\section{v1.4.0：两级报告与首末趋势}
\subsection{同一序列的双面板}
\key{\gradsparkcompare} 用第一位置和最后位置生成左侧哑铃，右侧保留全过程趋势。
两面板共用一个明确的范围和单位，要求固定范围或命名尺度，且仅支持线型趋势。
\begin{example}[text and listing,sidebyside=false]{项目得分：首末缺失不跳过}
\begingroup
\gradscale{project-score}{min=0,max=100}
\gradbarssetup{width=65mm,height=5ex,
 precision=0,label width=20mm,compare label=delta}
\begin{tabular}{lc}
项目 & 首末对比与全过程趋势\\
\hline
Atlas & \gradsparkcompare[scale group=project-score]
 {20,40,NA,80}\\
Birch & \gradsparkcompare[scale group=project-score]
 {NA,30,60,90}\\
Elm & \gradsparkcompare[scale group=project-score]
 {20,50,65,NA}\\
\end{tabular}
\endgroup
\end{example}
NA 保留位置；缺失端点不画标记及连接线，有效另一端仍保留，但差值缺失。
单点同时作为首末，空序列报错。默认 \key{compare label=values} 显示首末值，
\key{delta} 显示同单位的末减首；显式 text 覆盖标签，不支持 percent 对比标签。

\key{width} 包含双面板、默认 3mm 间距和两端各“标记尺寸加 1pt”的外延余量。
\key{compare panel ratio=.4} 分配剩余图宽，须严格在 0 与 1 之间，间距非负。
每个面板须宽于两倍余量；外置标签另占空间。一次组合生成一个 Figure，描述完整序列、
首末与差值，继续遵循显式 alt、artifact/off 和嵌入 TikZ 的既有规则。
\clearpage
\subsection{显式 CSV 两级文本分组}
\key{group levels} 接收恰好两个不同的已声明文本列 ID，与旧 \key{group by} 及手工组分离，不能混用。
演示 CSV 交错排列 North 与 South，两者均含名为 Review 的小组。
\begin{example}[text and listing,sidebyside=false,breakable]{区域项目评审：可见均值与父级排名}
\gradbarsloadcsv{projects}{docs/data/hierarchical-teams.csv}
\begin{center}
\begin{gradtable}[data=projects,style=hierarchical,
 width=\linewidth,group levels={region,team},
 rank scope=parent,parent summary=true,leaf summary=true,
 parent summary label={区域均值},
 leaf summary label={小组均值},label width=17mm]
\gradcolumn[source=Project,width=22mm]{project}{项目}
\gradcolumn[source=Region,width=17mm]{region}{区域}
\gradcolumn[source=Team,width=18mm]{team}{小组}
\gradcolumn[source=Score,type=number,width=20mm,
 aggregate=mean,mark=both,options={precision=1}]{score}{得分}
\gradcolumn[source=Trend,type=trendcompare,
 options={min=0,max=100,height=5mm,precision=0,
 compare label=delta}]{trend}{首末 / 趋势}
\gradsummary{可见行均值}
\end{gradtable}
\end{center}
\end{example}
父组和叶组按原始 CSV 首次出现顺序组织。结构化路径保留父组身份与子组原文，
不同父组下的同名子组不合并；裁剪首尾空格，空键使用缺失组标题，而字面 NA 是独立文本键。
逗号、斜杠不是路径分隔符。
\subsection{选择顺序、作用域与原生复合列}
流程继续为筛选、全局稳定排序、全局 limit、分组；叶组内部保留已选排序，组保持原始首次出现顺序。
不增加隐式组内 limit 或组内尺度。

\key{rank scope=table} 仍为默认；\key{parent} 比较同父组全部可见叶组的数据，
\key{leaf} 只比较同一路径。按规范原值判断并列和第二不同值，缺失不参与。
两级分组不使用旧 \key{rank scope=group}，摘要不参与排名和尺度。

\key{parent summary} 与 \key{leaf summary} 默认 false，可分别启用；各自拥有 label 与 scope 键。
\key{parent summary scope}、\key{leaf summary scope} 接受 visible 或 data，默认 visible。
data 使用该节点全部绑定原始行，可含被筛选或 limit 移除的行；完全不可见组不显示标题或摘要。
列 aggregate 仍须明确声明，总 \key{\gradsummary} 拥有独立作用域；旧单层摘要键不能混用。

原生 \key{type=trendcompare} 接收规范数字列表或 CSV 引号内序列，默认高 4ex，双面板一次解析同一尺度。
\key{scale=column} 读取全部有效观测并遵循表格尺度作用域，不把复合序列作为标量筛选、排序、
标记、聚合或地区输入。没有新增任意自定义列适配或 tabularray 辅助命令。

\subsection{长表与排版边界}
\key{style=hierarchical} 设置无障碍配色、行距与缩进叶标题。
原生 \key{long=true} 中父子标题连接首行，逐级摘要链连接相应末行；数据行之间可以分页。
单行或完整连接链须能放入一页。booktabs、grid、rows、none 四种横线样式均以实际 PDF 出页页码验证。
本版不新增完整层级 Tagged PDF 表结构保证或 PDF/UA 认证。

完整 72 行跨页报告见 \texttt{docs/examples/hierarchical-longtable.tex}，需编译两次。
中英文短例位于 \texttt{docs/examples/}；契约见 \texttt{docs/hierarchical-tables.md} 与
\texttt{docs/trend-compositions.md}。自定义标量后端须为 \key{spark-compare} 显式启用 TikZ 回退，后端模式仍为 1。
